\section{Conclusions}
\label{sec:conclusions}

This work shows that dense visual retrieval is the dominant signal for multimodal search over egocentric lifelogging video. \textbf{SigLIP2} achieves a mean Precision@10 of 0.620, substantially outperforming the BM25-only baseline at 0.070. The best overall effectiveness is obtained by Approach C, the \fusionname pipeline, which reaches a mean Precision@10 of 0.640. Its gains on Q1 and Q9, both improving by 0.1 over the visual baseline, show that transcript and caption evidence can provide complementary context for activity-oriented and symbol-recognition queries. Cross-encoder re-ranking with the \textbf{bge-reranker-v2-m3} model over the top-50 candidates further improves precision beyond the bi-encoder stage alone.

Several components of the system behaved as intended. The OCR hallucination gate removed 178,914 ungrounded \textbf{Florence-2} OCR-bearing caption records at a 100\% strip rate, preventing systematic noise from entering both the sparse BM25 index and the dense caption embedding space. Temporal deduplication enforced result diversity by suppressing frames within a $\pm 10$ second window of stronger results from the same stream, while cross-stream flagging exposed valid multi-angle evidence. The interactive Rocchio-style refinement mechanism also supports user feedback without rebuilding indexes.

The study also exposed clear limitations. BM25 retrieval is not sufficient as a standalone method for this domain because relevant visual content is rarely described verbatim in co-occurring ASR transcripts. Query Q7 reaches only 0.3 under both Approaches B and C because global frame-level embeddings do not resolve instance-level object identity well in cluttered scenes. Query Q10 remains at 0.1 under both approaches, suggesting that transient physical state is difficult to capture with static frame-level similarity alone.

Future work should therefore focus on richer localisation and temporal modelling. Region-level grounding methods such as \textbf{GroundingDINO} \cite{liu2024groundingdino} could help address instance-level failures such as Q7. Larger re-ranking pools may further improve precision for difficult queries. Temporal sequence modelling, as explored in \cite{lei2021less}, could support state-aware retrieval for transient object conditions such as Q10. A final direction is to replace hand-tuned fusion weights with adaptive weighting learned from relevance feedback.
